\section{The shared C1--C2 boundary and a positive source partition}
\label{sec:r13-source-overlap}

This section records an exact source-only identity. It does not by itself
admit any new independent testing domain. Its purpose is to identify the
same physical boundary in the two remaining source charts and to retain
the nonnegative slack when that boundary opens.

Use the original weak-C2 variables $B=H+L$, $z,y>0$ and $0<r<1$,
with the unchanged source gaps and nonnegative $n$. Write $e=1-r$,
$K=1+z$, $T_c=B+zr$, and let $\delta$ be the original positive
C2 denominator. Define the positive source polynomial
\[
 P_n=B+B^2+2Bzr+BHz+yB(1+z)+Hrz^2+yz^2re.
 \tag{OP1}
\]
Then the following identity holds on the entire original source:
\[
 \boxed{\quad eP_n=(1+z)(H+y)(B+zr)+2r\delta n.\quad}
 \tag{OP2}
\]
Every factor in the first summand is positive, and the second summand
is nonnegative, including literal $n=0$. In particular the pure-source
weights
\[
 q={ (1+z)(H+y)(B+zr)\over eP_n}>0,
 \qquad s_n={2r\delta n\over eP_n}\ge0
 \quad\hbox{satisfy}\quad q+s_n=1.                   \tag{OP3}
\]
Both square roots are legitimate source-field elements. Neither weight
is assumed to have nonzero residue.

\begin{proof}
The unchanged formula for $n$ is
\[
 2r\delta n=(B+zr)(1+zr)(1+zr+B-H)-r\delta.
\]
Substitution of the printed polynomial for $\delta$, followed by
collecting terms, gives OP2. This is a polynomial identity, so no
source boundary is obtained by continuity and no opening has been
inverted except the displayed positive pure-source factors.
\end{proof}

\paragraph{Exact correspondence of the source contacts.}
Use a subscript $1$ for native C1 notation in this paragraph. The map
\[
 x_1=z,\qquad \eta_1={y\over B},\qquad
 \zeta_1={B\over z},\qquad r_1={H\over B},\qquad
 \theta_1=e,
 \tag{OP4}
\]
with $k_1=\sqrt z$ and the native parameter
$\Delta_1=B/(z\sqrt y)$, gives
\[
 (u_c,v_c;W_c,X_c,Y_c,Z_c)
 =\left(\sqrt z,\sqrt y;{B\over\sqrt{zy}},
                       -r\sqrt z,{H\over\sqrt y},1\right)
\]
from the original native C1 contact formula. Thus it is the same
contact, including all state coordinates. Direct substitution in the
native positive polynomial $P_d$ gives $BP_d=P_n$.
The native C1 beta coefficient at this formal contact therefore obeys
\[
 \beta_1=-{r\delta n\over eP_n}=-{s_n\over2},
 \qquad \beta_1+\tfrac12={q\over2}>0.                 \tag{OP5}
\]
For $n=0$, the original C2 parent has
$\lambda=n=\beta=0$. It is literally an original C1 parent,
up to an optional positive normalization of $C$. The native source
entry and uniqueness at stationary contact identify its coefficients
with OP4. For $n>0$, OP5 makes the formal C1 beta negative. Hence
OP4 does \emph{not} map that source to the nonnegative C1 event cone,
and no physical C1 all-tests theorem may be imported there.

\paragraph{The revision-12 small-gap chart.}
Put $z=aB$, $y=tB$, $H=hB$, and $B=e\ell$, with the exact
BS1/GR1 hypotheses. Here $a,t$ are positive bounded units,
$h,e$ are infinitesimal, $t_0>1$, and $\ell$ is positive bounded.
Canceling the pure-source factor $e$ before taking residues gives
\[
 q_0=\ell_0t_0(1+a_0),\qquad
 (s_n)_0=1-\ell_0t_0(1+a_0).                          \tag{OP6}
\]
This is consistent with the original constraint $n\ge0$ and keeps
its equality case. At $n=0$, OP4 has
\[
 L_1=x_1\zeta_1=B,\quad
 \eta_1=t,\quad\zeta_1=1/a,\quad r_1=h,\quad
 (j_{0,1})_0=t_0-1>0.
\]
It therefore enters the weighted small-$L$ C1 theorem in
Section~\ref{sec:r13-c1-r-openings}. This overlap is an exact reuse
of a proved physical boundary, not a claim that its positive-$n$
neighborhood is covered. All total-scale and earlier-complement
intersections remain in force.
